\scanpage{96}
\uline{Prop.} 3 \quad Si $t$ est un triangle contenant une arête Sup. transverse (i.e. contient un pt d'un $\widetilde E_\alpha$ et un d'un $\widetilde E_\beta$, avec \inserted{$\alpha,\beta\in t_0$} $\alpha\ne\beta$) alors les trois arêtes sont transverses, et la relation entre $\alpha\in t_0$ et $s\in t$ : ``$\alpha$ est lié à $s$ (i.e. $s\in\widetilde E_\alpha$)'' est une bijection $t_0\xrightarrow{\sim}t$.

Les triangles en question s'appellent triangles supérieurs, leur \unsure{2} ens est noté $\deleted{\unsure{\widetilde T}}\widetilde\Gamma$. On a donc des applications
\[
\deleted{\unsure{\widetilde T}}\widetilde\Gamma\xrightarrow{\widetilde\varphi_\alpha}\widetilde E_\alpha\qquad(\alpha\in t)
\]
(ou encore
\diagram{\begin{tikzcd}[ampersand replacement=\&]
\overline E_a\& {}\&\widetilde E_b\\
{}\&\widetilde\Gamma\arrow[ul,"\overline\varphi_a"]\arrow[ur,dash,"\widetilde\varphi_b"']\arrow[d,"\widetilde\varphi_c"]\&{}\\
{}\&\widetilde E_c\&{}
\end{tikzcd}}
) \deleted{définissant}\inserted{\unsure{équivaut à|équivalant à|équivalentes à}} une application unique
\[
\widetilde\Gamma\xrightarrow{\widetilde\varphi}\prod_{\alpha\in t_0}\widetilde E_\alpha
\qquad(=\widetilde E_c\times\widetilde E_b\times\overline E_c),
\]
L'axiome 1 pour les arêtes sup. transverses équivaut alors à la

\uline{Prop.} 4 \quad Si $\deleted{\alpha,\beta\in t_0,\ \alpha=}u\in\unsure{\mathcal P|\mathfrak P}_2(t_0)$ est une
% archive sheet number: 95
